\documentclass[10pt]{article}

\usepackage[utf8]{inputenc}

\usepackage[T1]{fontenc}

\usepackage{amsmath}

\usepackage{amsfonts}

\usepackage{amssymb}

\usepackage[version=4]{mhchem}

\usepackage{stmaryrd}

\usepackage{bbold}



\begin{document}

М. м. (то же самое справедливо и для каскадов; о более общих группах преобразований см., напр., [3] или [4]). Сформулированное свойство точки $w$ (и ее траектории) Дж. Бирктоф назвал рекуррентн о сть ю этой точки (и траектории), употребителен также предложенный У. Готшалком и Г. Хедлундом [3] другой вариант терминологии, в к-ром это свойство называется почт и п ер и одич н ос тью точки $w$. Если $F=W$, то саму динамич. систему называют м и н и мал ь н о й.



Если нек-рая траектория имеет компактное замыкание, то в нем содержится нек-рое М. м. F (для полугрупп непрерывных преобразований $\left\{S_{t}\right\}$ с неотрицательными действительными или целыми $t$ справедлив аналог этого результата, причем на $F$ преобразования $S_{t}$ уже обратимы [5]). Однако исследование предельного поведения траекторий динамич. системы не сводится к изучению одних только М. М. последней. Очень просто устроено М. м. гладкого потока класca $C^{2}$ на двумерной замкнутой поверхности - это либо точка, либо замкнутая траектория, либо вся поверхность, к-рая в этом случае является тором (т е о ре ма Ш в а р ца, см. [6]). В общем случае строение М. м. может быть весьма сложным (в связи с этим, дополнительно к сказанному в [2]-[4], надо указать, что минимальность динамич. системы не накладывает никаких ограничений на ее эргодич. свойства по отношению к имеющимся у нее инвариантным мерам [7]). М. м.- основной объект изучения в топологич. динамике.



Лит.: [1] Б и р к г о ф Дж. Д., Динамические системы, пер. с англ., М.-Л., 1941; [2] Немыцки й В. В., Сте пан ов В. В., Качественная теория дифференциальных уравнений, 2 изд., М.- Л., 1949; [3] Gottschalk W. H., Hedlund G. A., Topological dynamics, Providence, 1955; [4] Бр о нш те й н И. У.,' Расширения минимальных групп преобразований, Киш., 1975; [5] Л е в и тан Б. М., Ж и ко в В. В., Почти-периодические функции и дифференщиальные уравнения, М., 1978; [6] Х а р т а н Ф., Обыкновенные дифференциальные уравнения, пер. с англ., М., 1970; [7] Итоги науки и техники. Математический анализ, т. 13, М., 1975, с. 129-262.



МИНИМИЗИРУКЩАЯ ПОСЛЕДОВАТЕЛЬНОСТЬ см. Дирихле принцип.



МИНКОВСКОГО ПРОСТРАНСТВО-ВРЕМЯ - четЫрехмерное псевдоевклидово пространство сигнатуры $(1,3)$, предложенное Г. Минковским (H. Minkowski, 1908) в качестве геометрической интерпретации пространства-времени специальной теории относительности. «Воззрения на пространство-время... возникли на экспериментально-физической основе. В этом их сила. Их тенденция радикальна. Отныне пространство само по себе и время само по себе должны обратиться в фикции, и лишь некоторый вид соединения обоих должен сохранить самостоятельность» (cM. [1]).



Каждому событию соответствует точка М. п.-в., координаты трехмерного физич. пространства, а четвертая - координата $c t$, где $c$ - скорость света, $t$ - время события. Обычно используются прямоугольные декартовы координаты $x^{1}$, $x^{2}, x^{3}$ события в нек-рой инерциальной системе отсчета, $x^{0}=c t$. Геометрич. свойства М. п.-в. определяются выражением для квадрата так наз. и нтер вала между двумя событиями



\[

s^{2}=\left(\Delta x^{0}\right)^{2}-\sum_{i}\left(\Delta x^{i}\right)^{2} .

\]



При переходе от одной инерциальной системы отсчета к другой пространственные координаты и время преобразуются друг в друга посредством Пуанкаре преобразований. Геометрия М. П.-в. позволяет наглядно интерпретировать кинематич. релятивистские эффекты (сокращение длины, замедление времени и т. д.).



К открытию фундаментальной важности о том, что пространство и время едины и геометрия его псевдоевклидова, привело изучение электромагнитных явлений (хотя даже тогда, когда были открыты уравнения Максвелла-Лоренца, никто не предполагал, что они изменят сложившееся представление о пространстве и времени. Этим открытием и обязаны мы в высшей степени Г. Минковскому, несмотря на то, что главное в этом направлении уже содержалось в работах А. Пуанкаре (Н. Poincaré) (см. [1]). Именно он первый установил, что величина, впоследствии названная интервалом, инвариантна относительно Лоренца преобразований.



Лит.: [1] Принцип относительности, сб. статей, М., 1973; [2] Л о г у н о в А. А., Лекции по теории относительности и гравитации, М., 1987; [3] е го ж е, К работам Анри Пуанкаре «О динамике электрона», М., $1984 . \quad$ М. И. Войцеховский. МИРОВАЯ ЛИНИЯ - линия в пространстве-времени, являющаяся траекторией движущейся точки или последовательностью событий. Впрочем, вернее было бы говорить об изображении не «процесса движения», а об «истории существования» данной точки. Дело в том, что движение рассматривалось всегда относительно той или иной системы отсчета, между тем как М. л. является построением абсолютным, не зависящим от выбора системы отсчета и в таком выборе вообще не нуждающимся (см. [1]).



При этом под с обы т и м понимается элементарное событие, то есть происходящее в столь малой области пространства и в столь короткий промежуток времени, что, идеализируя положение вещей, их можно считать происходящими в одном месте (с пространственными координатами $(x, y, z)$ ) и мгновенно (в момент времени $t$ ).



Обычно рассматриваются гладкие (или кусочно гладкие) М. л.; они разделяются на три класса: пространственноподобные, изотропные и времениподобные в зависимости от того, каким является четырехмерный вектор, касательный к М. л. Так, напр., М. л. материальной точки с положительной массой покоя является пространственноподобной кривой; М. л. материальной точки с нулевой массой покоя (такая точка является неквантовой моделью фотона, нейтрино и других элементарных частиц нулевой массы) является изотропной кривой.



Лит.: [1] Р а ш в с к и й П. К., Риманова геометрия и тензорный анализ, М., 1964



М. И. Войцеховский.



МИРОВАЯ ТОЧКА - точка пространства-времени, описывающая событие, происшедшее в данной точке пространства в некоторый момент времени. М. И. Войцеховский. МИTТАГ-ЛЕФФЛЕРА TEOPEMA $о$ разложен и и мероморфной функци и - одна из основных теорем теории аналитических функций, дающая для мероморфных функций аналог разложения рациональной функщии на простейшие дроби. Пусть $\left\{a_{n}\right\}_{n=1}^{\infty}$ - последовательность различных комплексных чисел



\[

\left|a_{1}\right| \leqslant\left|a_{2}\right| \leqslant \ldots, \quad \lim _{n \rightarrow \infty} a_{n}=\infty,

\]



и $\left\{g_{n}(z)\right\}$ - последовательность рациональных функций вида:



\[

g_{u}(z)=\sum_{k=1}^{l_{n}} \frac{c_{n k}}{\left(z-a_{n}\right)^{k}},

\]



так что точка $a_{n}$ является единственным полюсом соответствующей функции $g_{n}(z)$. Тогда существуют мероморфные функции $f(z)$ в плоскости $\mathbb{C}$ комплексного переменного $z$, имеющие полюсы в точках $a_{n}$ и только в них с заданными главными частями (1) Лорана рядов, соответствующих точкам $a_{n}$. Все эти функции $f(z)$ представимы в виде р а з ло жения Миттаг-Леффлера:



\[

f(z)=h(z)+\sum_{n=1}^{\infty}\left|g_{n}(z)+p_{n}(z)\right|,

\]



где $p_{n}(z)$ - нек-рые многочлены, подбираемые по $a_{n}$ и $g_{n}(z)$ так, чтобы ряд (2) равномерно (после выбрасывания конечного числа членов) сходился на любом компакте $K \subset \mathbb{C}, h(z)$ - произвольная целая функция.



Из М.-Л. т. вытекает, что любая наперед заданная мероморфная функция $f(z)$ в $\mathbb{C}$ с полюсами $a_{n}$ и соответствуюшими главными частями $g_{n}(z)$ разложений $f(z)$ в ряды Лорана в окрестности $a_{n}$ разлагается в-ряд (2), где целая





\end{document}